\subsection{Dualność Pontrjagina}

Cel: $G\cong \Hat{\Hat{G}}$ przez $x\mapsto \langle x, -\rangle$.

\begin{theorem}{Bochnera}{}
  Jeśli $\phi\in\mathcal{P}(G)$ to istnieje jedyna dodatnia miara $\mu\in M(\Hat{G})$ taka, że $\phi=\phi_\mu$.

  \color{red}Gal robi dla niedodatnich tylko tak ogólniej?
\end{theorem}

\begin{proof}
  BSO $\phi(1)=1$. Definiujemy formę Hermitowską
  $$\langle f,g\rangle_\phi=\int\phi(g^*\ast f)$$
  z nierówności Schwarza:
  $$\left|\int\phi(g^*\ast f)\right|^2\leq \int \phi(f^*\ast f)\int \phi(g^*\ast g)$$
  dla $f,g\in L^1G$.
\end{proof}

\begin{zadanie}{3}{10}
  Niech $f\in L^1A\cap L^2A$. Niech $\Hat{\eta}$ będzie współczynnikiem diagonalnym $f$ (z twierdzenia Bochnera). Zrozum napis (wiadomo co to całka z miary; całka z funkcji to całka względem miary Haara):
  $$\|f\|_2^2=\taH{\eta}(1)=\int\eta=\int\Hat{\taH{\eta}}=\int|\Hat{f}|^2=\|\Hat{f}\|^2_2.$$
\end{zadanie}

\begin{proof}
  Co robi Folland? Folland mówi
  $$\|f\|_2^2=\int f^2=(f\ast f^*)(1)=\int\Hat{f\ast f^*}(\xi)d\xi=\int |\Hat{f}(\xi)|^2d\xi,$$
  gdzie 
  $$f\ast f^*(x)=\int f(y)f^*(y^{-1}x)dy=\int f(y)\Delta(x^{-1}y)\overline{f(x^{-1}y)}=\int f(y)\overline{f(x^{-1}y)}$$
  bo mamy grupę abelową.

  $\color{green}\|f\|_2^2=\taH{\eta}(1)$

  Moje podejrzenia: bierzemy $f$ jak wyżej i robimy z tego na siłę funkcję dodatnio określoną, czyli $f\ast f^*$. Z twierdzenia Bochnera mamy $\eta\in M(\Hat{G})$ takie, że 
  $$(f\ast f^*)(x)=\int\langle x, \xi\rangle d\eta(\xi)(=\Hat{\eta}(x)).$$

  Czyli jak u Follanda dla miary $\mu\in M(\Hat{G})$ robimy funkcję z $C_b(G)$ przez $\mu\mapsto \phi_\mu$ i odwracamy to przez $\phi\mapsto \mu_\phi$, to u Gala to odwracanie to jest $\taH{\phi}=\mu_\phi$. 
  
  $\color{green}\taH{\eta}(1)=\int \eta$

  Z twierdzenia Fouriera o inwersji/odwracaniu (?) mam
  $$\taH{\eta}(x)=\int \langle x, \xi\rangle d\eta(\xi)$$
  no a jeśli $x=1$ to mam po prostu po prawej $\int d\eta(\xi)$.

  $\color{green}\int\eta=\int\Hat{\taH{\eta}}$

  Z twierdzenia Fouriera o odwracaniu (?) wiemy, że $d\eta(\xi)=\Hat{\taH{\eta}}(\xi)d\xi$. 

  $\color{green}\int\Hat{\taH{\eta}}=\int |\Hat{f}|^2$
  
  Po pierwsze z definicji $f^*$ i transformaty Foueriera mam
  \begin{align*}
    \Hat{f^*}(\xi)&=\int \overline{\langle x, \xi\rangle}f^*(x)dx=\int\overline{\langle x,\xi\rangle f(x^{-1})}dx=\\ 
                  &=[x^{-1}=u,\;dx=du]=\int \overline{\langle u^{-1}, \xi\rangle f(u)}du=\\ 
                  &=\int\langle u,\xi\rangle \overline{f(u)}du=\overline{\int \overline{\langle x, \xi\rangle}f(x)dx}=\\ 
                  &=\overline{\Hat{f}(\xi)}
  \end{align*}

  Tutaj wiem, że $\taH{\eta}=f\ast f^*$, czyli
  $$\int \Hat{\taH{\eta}}(\xi)d\xi=\int \Hat{f\ast f^*}(\xi)d\xi=\int \Hat{f}(\xi)\Hat{f^*}(\xi)d\xi$$
  i z wyliczenia wyżej dostaję $\int \Hat{f}(\xi)\overline{\Hat{f}(\xi)}d\xi=\int |\Hat{f}(\xi)|^2d\xi$.
\end{proof}

\begin{zadanie}{3}{11}
  Wywnioskuj twierdzenie Reigleya-Parsevala-Plancherela: transformata Fouriera jednoznacznie rozszerza się do unitarnej izometrii $L^2A\to L^2\Hat{A}$.
\end{zadanie}

\begin{proof}
  Zbiór $L^1 G\cap L^2G$ jest gęsty w $L^2G$, bo Stone-Weierstrass lub pewnie można delikatniej. 

  Zadanie wcześniej pokazuje, że $f\mapsto \Hat{f}$ jest izometrią w $L^2$ normie na $L^1G\cap L^2 G$. Czyli rozszerza się do izometrii $L^2G\to L^2\Hat{G}$ (która zachowuje granice ciągów -> chyba w ten sposób bym ja konstruowała wogóle).

  Czyli mamy w ręku izometrię $F:L^2G\to L^2\Hat{G}$, czyli będzie $F^*F=Id$. Teraz że $FF^*=Id$, czyli że jest unitarna -> Folland pokazuje surjektywność. Tutaj mamy elegancki iloczyn skalarny, czyli żeby pokazać surjektywność $F$ wystarczy pokazać, że $F^\perp=0$. W takim razie bierzemy $\phi$ ortogonalne do wszystkich $\Hat{f}$ dla $f\in L^1G\cap L^2G$. To oznacza, że 
  $$0=\langle \phi, \Hat{f}\rangle_{L^2}=\int \phi\overline{\Hat{f}}$$
  w szczególności gdy pod $f$ podstawimy $L_\omega f$ dla dowolnego $\omega$ to dostajemy
  $$0=\int \phi\overline{\Hat{L_\omega f}}=\int \phi(\xi)\overline{\overline{\langle\omega, \xi\rangle}\Hat{f}(\xi)}=\int \phi(\xi)\langle \omega, \xi\rangle \overline{\Hat{f}(\xi)}.$$
  Dalej, skoro $\phi, \Hat{f}\in L^2\Hat{G}$, to $\phi\Hat{f}\in L^1\Hat{G}$, czyli $\phi(\xi)\overline{\Hat{f}(\xi)}d\xi$ definiuje miarę na $\Hat{G}$. Ale $\mu\to \phi_\mu$ dla $\phi_\mu(x)=\int \langle x, \xi\rangle d\mu(\xi)$ jest injekcją, więc skoro tutaj dla dowolnego $f$ mamy zero, to znaczy że miara $\phi\overline{\Hat{f}}d\xi$ jest felerna, czyli $\phi\overline{\Hat{f}}=0$ dla wszystkich $f$, czyli $\phi=0$.
\end{proof}

